\setcounterpageref{ex}{0}
\TaoPhieuKiemTra{06}

\setcounter{ex}{0}
\setcounter{bt}{0}

\textbf{\large I. PHẦN TRẮC NGHIỆM (3,0 điểm)}

{\itshape Khoanh tròn vào chữ cái đặt trước câu trả lời đúng. Mỗi câu đúng được 0,5 điểm.}

\begin{ex} % Câu 1
Số $75$ đọc là
\choice
{Bảy mươi năm}
{Bảy lăm}
{\True Bảy mươi lăm}
{Năm mươi bảy}
\loigiai{
Số $75$ đọc là Bảy mươi lăm.\\
\textbf{Đáp án đúng: C}
}
\end{ex}

\begin{ex} % Câu 2
Tổng của hai số hạng $30$ và $40$ là
\choice
{$60$}
{\True $70$}
{$80$}
{$10$}
\loigiai{
Tổng: $30 + 40 = 70$.\\
\textbf{Đáp án đúng: B}
}
\end{ex}

\begin{ex} % Câu 3
Tính: $15\text{ kg} - 7\text{ kg} = \dots$
\choice
{$7\text{ kg}$}
{$9\text{ kg}$}
{$6\text{ kg}$}
{\True $8\text{ kg}$}
\loigiai{
$15 - 7 = 8\text{ kg}$.\\
\textbf{Đáp án đúng: D}
}
\end{ex}

\begin{ex} % Câu 4
Hình tứ giác có tất cả bao nhiêu cạnh?
\choice
{\True $4$ cạnh}
{$3$ cạnh}
{$5$ cạnh}
{$6$ cạnh}
\loigiai{
Hình tứ giác có $4$ cạnh và $4$ đỉnh.\\
\textbf{Đáp án đúng: A}
}
\end{ex}

\begin{ex} % Câu 5
Điền dấu thích hợp: $16 - 8 \dots 17 - 9$.
\choice
{$>$}
{\True $=$}
{$<$}
{Không so sánh được}
\loigiai{
$16 - 8 = 8$; $17 - 9 = 8$. Vì $8 = 8$ nên điền dấu $=$.\\
\textbf{Đáp án đúng: B}
}
\end{ex}

\begin{ex} % Câu 6
Tính: $15 - 9 + 6 = \dots$
\choice
{$11$}
{$13$}
{\True $12$}
{$10$}
\loigiai{
$15 - 9 = 6$; $6 + 6 = 12$.\\
\textbf{Đáp án đúng: C}
}
\end{ex}

\vspace{4mm}
\textbf{\large II. PHẦN TỰ LUẬN (7,0 điểm)}

{\itshape Thí sinh trình bày chi tiết bài làm vào giấy kiểm tra.}

\begin{bt}[1,0 điểm]
Tính nhẩm:
\begin{multicols}{2}
\begin{enumerate}[a)]
  \item $8 + 5$
  \item $7 + 6$
  \item $13 - 7$
  \item $15 - 8$
\end{enumerate}
\end{multicols}
\loigiai{
\begin{itemize}
  \item $8 + 5 = 13$
  \item $7 + 6 = 13$
  \item $13 - 7 = 6$
  \item $15 - 8 = 7$
\end{itemize}
}
\end{bt}

\begin{bt}[1,5 điểm]
Đặt tính rồi tính:
\begin{multicols}{3}
\begin{enumerate}[a)]
  \item $35 + 50$
  \item $87 - 39$
  \item $53 + 44$
\end{enumerate}
\end{multicols}
\loigiai{
\begin{itemize}
  \item $35 + 50 = 85$
  \item $87 - 39 = 48$
  \item $53 + 44 = 97$
\end{itemize}
}
\end{bt}

\begin{bt}[1,0 điểm]
Tính giá trị biểu thức kèm đơn vị đo:
\begin{enumerate}[a)]
  \item $22\text{ kg} + 9\text{ kg} - 9\text{ kg}$
  \item $31\text{ }l - 12\text{ }l + 8\text{ }l$
\end{enumerate}
\loigiai{
\begin{enumerate}[a)]
  \item $22\text{ kg} + 9\text{ kg} - 9\text{ kg} = 31\text{ kg} - 9\text{ kg} = 22\text{ kg}$.
  \item $31\text{ }l - 12\text{ }l + 8\text{ }l = 19\text{ }l + 8\text{ }l = 27\text{ }l$.
\end{enumerate}
}
\end{bt}

\begin{bt}[1,5 điểm]
Bạn Khánh có $21$ cây bút chì. Bạn Linh có ít hơn bạn Khánh $4$ cây bút chì. Hỏi bạn Linh có bao nhiêu cây bút chì?
\loigiai{
Số cây bút chì bạn Linh có là:
\[ 21 - 4 = 17 \text{ (cây bút chì)} \]
\textbf{Đáp số:} $17$ cây bút chì.
}
\end{bt}

\begin{bt}[1,0 điểm]
Một thùng chứa đầy $41\text{ }\text{kg}$ gạo nếp. Buổi sáng cửa hàng đã bán ra $18\text{ }\text{kg}$. Hỏi trong thùng còn lại bao nhiêu ki-lô-gam gạo nếp?
\loigiai{
Trong thùng còn lại số ki-lô-gam gạo nếp là:
\[ 41 - 18 = 23\text{ }\text{kg} \]
\textbf{Đáp số:} $23\text{ }\text{kg}$ gạo nếp.
}
\end{bt}

\begin{bt}[1,0 điểm]
Bạn Nam có một số viên bi. Nếu Nam cho Bình $5$ viên bi thì Nam còn lại $8$ viên bi. Hỏi lúc đầu Nam có bao nhiêu viên bi?
\loigiai{
Lúc đầu Nam có số viên bi là:
\[ 8 + 5 = 13 \text{ (viên bi)} \]
\textbf{Đáp số:} $13$ viên bi.
}
\end{bt}
